\subsection{好滤过}

设 A 为滤过交换环，E 为滤过 A-模，\((A_{n})\) 和 \((E_{n})\)

分别是 A 与 E 的滤过；设 \(A_{0} = A\)。在多项式环 A{[}X{]} 中，集合 \(A' = \sum_{n \geqslant 0} A_{n} X^{n}\) 是一个 N 型分次子 A-代数；子群 \(E' = \sum_{n \geqslant 0} E_{n} \otimes_{A} A X^{n}\)（属于 \(E \otimes_{A} A[X]\)）是一个 N 型分次 \(A'\)-模，因为

\[
\mathbf {A} _ {m} \mathbf {X} ^ {m} (\mathrm{E} _ {n} \otimes_ {\mathbf {A}} \mathbf {A X} ^ {n}) \subset (\mathbf {A} _ {m} \mathrm{E} _ {n} \otimes_ {\mathbf {A}} \mathbf {A X} ^ {m + n}) \subset \mathrm{E} _ {m + n} \otimes_ {\mathbf {A}} \mathbf {A X} ^ {m + n}.
\]

定义 1。设 A 为交换环，m 为 A 的理想，E 为 A-模，\((\mathrm{E}_{n})\) 是加法群 E 上由 E 的子模组成的一个滤过。若满足以下条件，则称滤过 \((\mathrm{E}_{n})\) 为 m-好：

\begin{enumerate}
\def\labelenumi{(\arabic{enumi})}
\item
  \(\mathfrak{m}\mathrm{E}_n\subset \mathrm{E}_{n + 1}\) 对所有 \(n\in \mathbf{Z}\) 都成立；
\item
  存在整数 \(n_0\)，使得 \(\mathfrak{m}\mathrm{E}_n = \mathrm{E}_{n+1}\) 对 \(n \geqslant n_0\) 成立。
\end{enumerate}

于是由 \(q\) 归纳可得，\(\mathfrak{m}^q\mathrm{E}_n = \mathrm{E}_{n + q}\) 对 \(n \geqslant n_0\)、\(q \geqslant 1\) 成立。注意，若 A 赋予 m-进滤过，条件（1）意味着滤过 \((\mathrm{E}_n)\) 与 E 上的 A-模结构相容。显然，在每个 A-模 E 上，m-进滤过都是 m-好的。若 A-模 E 上的一个滤过是 m-好的，则 E 的任意商模上的商滤过也是 m-好的。

定理 1。设 A 为交换环，m 为 A 的理想，E 为 A-模，\((\mathrm{E}_{n})\) 是加法群 E 上由有限生成的子 A-模组成的滤过。假设对所有 n 都有 \(mE_{n} \subset E_{n+1}\)。令 \(A'\) 为由 \(\sum_{n \geqslant 0} m^{n} X^{n}\) 组成的 A{[}X{]} 的分次子代数，令 \(E'\) 为分次 \(A'\)-模 \(\sum_{n \geqslant 0} E_{n} \otimes_{A} AX^{n}\)。下列两个条件等价：

\begin{enumerate}
\def\labelenumi{(\alph{enumi})}
\item
  滤过 \((\mathbf{E}_n)\) 是 m-好的。
\item
  E' 是有限生成的 A'-模。
\end{enumerate}

假设 \(mE_{n-1} = E_n\) 对 \(n > n_0 \geqslant 0\) 成立。对于 \(i \leqslant n_0\)，令 \((e_{ij})_{1 \leqslant j \leqslant r_i}\) 为 A-模 \(E_i\) 的一个有限生成系。由于 A-模 \(E_n \otimes_A X^n\) 由元素 \(e_{nj} \otimes X^n\)（\(0 \leqslant n \leqslant n_0\)）生成，且等于

\[
\mathfrak {m} ^ {n - n _ {0}} \mathrm{E} _ {n _ {0}} \otimes_ {\mathbf {A}} \mathrm{AX} ^ {n}
\]

当 \(n > n_0\) 时，A'-模 E' 由元素 \(e_{nj} \otimes \mathbf{X}^n\)（\(0 \leqslant n \leqslant n_0\) 且 \(1 \leqslant j \leqslant r_n\)）生成；所以它当然是有限生成的。

反之，若 \(\mathbf{E}'\) 是有限生成的 \(\mathbf{A}'\)-模，则它由形如 \(e_k \otimes \mathbf{X}^{n(k)}\) 的元素组成的有限族生成，其中 \(e_k \in \mathbf{E}_{n(k)}\)。令 \(n_0\) 为各整数 \(n(k)\) 中的最大者。则对于 \(n \geqslant n_0\) 及 \(f \in \mathbf{E}_n\)，

\[
f \otimes \mathbf {X} ^ {n} = \sum_ {k} t _ {k} (e _ {k} \otimes \mathbf {X} ^ {n (k)})
\]

其中 \(t_k \in \mathbf{A}'\)；必要时将 \(t_k\) 替换为其次数为 \(n - n(k)\) 的齐次分量后，可设 \(t_k = a_k \mathbf{X}^{n - n(k)}\)，其中 \(a_k \in \mathfrak{m}^{n - n(k)}\)。

由于唯一元素 \(X^{n}\) 构成 A-模 \(AX^{n}\) 的一个基，等式 \(f \otimes X^{n} = \left( \sum_{k} a_{k} e_{k} \right) \otimes X^{n}\) 蕴含 \(f = \sum_{k} a_{k} e_{k}\)。于是 \(E_{n} \subset m^{n-n_{0}} E_{n_{0}}\)；由于反向包含显然成立，

\[
\mathrm{E} _ {n} = \mathfrak {m} ^ {n - n _ {0}} \mathrm{E} _ {n _ {0}},
\]

所以 \(\mathbf{E}_n = m\mathbf{E}_{n - 1}\) 当 \(n > n_0\) 时成立

引理 1。设 A 为交换 Noether 环，m 为 A 的理想。则由 \(\mathbf{A}' = \sum_{n\geqslant 0}\mathfrak{m}^n\mathbf{X}^n\) 组成的 A{[}X{]} 的子环是 Noether 环。

A' 是由 mX 生成的 A-代数；由于 A 是 Noether 环，mX 是有限生成的 A-模，结论于是由第 2 节第 10 号定理 2 的推论 3 得出。

命题 1。设 A 为交换 Noether 环，m 为 A 的理想；赋予 A m-进滤过。设 E、F 是两个滤过 A-模，且 j: F → E 是与滤过相容的单同态。若 E 有限生成且其滤过是 m-好的，则 F 有限生成且其滤过也是 m-好的。

由于 F 同构于 E 的一个子模，又因 A 是 Noether 环且 E 有限生成，故 F 有限生成。设 \((\mathbf{E}_{n})\)、\((\mathbf{F}_{n})\) 分别为 E、F 上的滤过，它们由有限生成子模组成；沿用引理 1 的记号，置 \(E' = \sum_{n \geqslant 0} E_{n} \otimes_{A} AX^{n}\)、\(F' = \sum_{n \geqslant 0} F_{n} \otimes_{A} AX^{n}\)；由于按假设 \(F_{n}\) 同构于 \(E_{n}\) 的一个子模，我们看到 \(F'\) 同构于 \(E'\) 的一个子模。根据定理 1，\(E'\) 是有限生成的 \(A'\)-模，因而 \(F'\) 也是有限生成的；又由于 \(A'\) 是 Noether 环（引理 1），由定理 1 即得结论。

推论 1（Artin–Rees 引理）。设 A 为交换 Noether 环，m 为 A 的理想，E 为有限生成的 A-模，F 为 E 的子模。由 E 上的 m-进滤过在 F 上诱导的滤过是 m-好的。

换言之，存在整数 \(n_{0}\)，使得

\[
\mathfrak {m} ((\mathfrak {m} ^ {n} \mathrm{E}) \cap \mathrm{F}) = (\mathfrak {m} ^ {n + 1} \mathrm{E}) \cap \mathrm{F}\tag{1}
\]

对所有 \(n \geqslant n_0\) 都成立。

推论 2。设 A 为交换 Noether 环，a、b 为 A 的两个理想。存在整数 h \textgreater{} 0，使得 \(a^{h} \cap b \subset ab\)。

存在 \(n\)，使得 \(\mathfrak{a}^{n+1} \cap \mathfrak{b} = \mathfrak{a}(\mathfrak{a}^n \cap \mathfrak{b}) \subset \mathfrak{ab}\)；这是将推论 1 应用于 \(\mathbf{E} = \mathbf{A}, \mathbf{F} = \mathfrak{b}\) 的结果。

推论 3。设 A 为交换 Noether 环，m 为 A 的理想，x 为 A 中不是零因子的元素。存在整数 \(k > 0\)，使得对所有 \(n \geqslant k\)，关系 \(xy \in \mathfrak{m}^n\) 蕴含 \(y \in \mathfrak{m}^{n - k}\)。

将推论 1 应用于 \(\mathbf{E} = \mathbf{A}\)、\(\mathbf{F} = \mathbf{A}x\)，可知存在 \(k\)，使得对所有 \(n \geqslant k\)，\(\mathfrak{m}^n \cap \mathbf{A}x = \mathfrak{m}^{n - k}(\mathfrak{m}^k \cap \mathbf{A}x)\)。于是，若 \(xy \in \mathfrak{m}^n\)，

\[
x y \in \mathfrak {m} ^ {n} \cap A x \subset \mathfrak {m} ^ {n - k} x
\]

又因 \(x\) 不是零因子，我们得到 \(y \in \mathfrak{m}^{n - k}\)。

用传递子（第一章，第 2 节，第 10 号）的记号，推论 3 的结论写作

\[
\mathfrak {m} ^ {n}: \mathrm{A} x \subset \mathfrak {m} ^ {n - k}.\tag{2}
\]

推论 4。设 A 为交换 Noether 环，m 为 A 的理想，E 为有限生成的 A-模，\((\mathbf{E}_{n})\) 和 \((\mathbf{E}_{n}^{\prime})\) 是由 E 的子模组成的两个滤过。假设当 A 赋予 m-进滤过时，滤过 \((\mathbf{E}_{n})\) 和 \((\mathbf{E}_{n}^{\prime})\) 与 E 上的 A-模结构相容。若滤过 \((\mathbf{E}_{n})\) 是 m-好的，且有 \(E_{n}^{\prime} \subset E_{n}\) 对所有 \(n \in Z\) 都成立，则滤过 \((\mathbf{E}_{n}^{\prime})\) 是 m-好的。这是命题 1 的特殊情形。

引理 2。设 A、B 为两个交换 Noether 环，\(\phi: \mathrm{A} \to \mathrm{B}\) 为环同态，E 为有限生成的 A-模，F 为有限生成的 B-模。则 \(\operatorname{Hom}_{\mathbf{A}}(\mathrm{E}, \phi_*(\mathrm{F}))\) 是有限生成的 B-模。

根据假设，存在满射 A-同态 \(v: \mathbf{A}^n \to \mathbf{E}\)；因此，映射 \(u \mapsto u \circ v\) 从 \(\operatorname{Hom}_{\mathbf{A}}(\mathbf{E}, \phi_*(\mathbf{F}))\) 到 \(\operatorname{Hom}_{\mathbf{A}}(\mathbf{A}^n, \phi_*(\mathbf{F}))\) 是单射；又因 B 是 Noether 环，只需证明 \(\operatorname{Hom}_{\mathbf{A}}(\mathbf{A}^n, \phi_*(\mathbf{F}))\) 是有限生成的 B-模即可；这是立即的，因为它同构于 \(\mathbf{F}^n\)。

命题 2。设 A 为交换 Noether 环，m 为 A 的理想，E、F 为两个有限生成的 A-模。若 \((\mathbf{F}_{n})\) 是 F 上的 m-好滤过，则子模 \(\operatorname{Hom}_{\mathbf{A}}(\mathbf{E}, \mathbf{F}_{n})\) 构成 A-模 \(\operatorname{Hom}_{\mathbf{A}}(\mathbf{E}, \mathbf{F})\) 上的 m-好滤过。由于 \(m^{k}F_{n}\subset F_{n+k}\) 对 \(n\in Z, k\geqslant0\) 成立，还成立：

\[
\mathfrak {m} ^ {k} \operatorname{Hom} _ {\mathbf {A}} (\mathrm{E}, \mathrm{F} _ {n}) \subset \operatorname{Hom} _ {\mathbf {A}} (\mathrm{E}, \mathrm{F} _ {n + k});
\]

于是族 \((\mathrm{Hom}_{\mathbf{A}}(\mathrm{E},\mathrm{F}_{n}))_{n\in\mathbf{Z}}\) 是 \(\mathrm{Hom}_{\mathbf{A}}(\mathrm{E},\mathbf{F})\) 上的滤过，并且与其作为赋予 m-进滤过的环 A 上的模结构相容。由于 E 有限生成，存在整数 r\textgreater0 及满射 A-同态 \(u:A^{r}\rightarrow E\)，它定义了一个单射 A-同态

\[
v = \operatorname{Hom} (u, 1 _ {\mathrm{F}}): \operatorname{Hom} _ {\mathrm{A}} (\mathrm{E}, \mathrm{F}) \rightarrow \operatorname{Hom} _ {\mathrm{A}} (\mathrm{A} ^ {r}, \mathrm{F});
\]

显然，\(v\) 与滤过 \((\mathrm{Hom}_{\mathbf{A}}(\mathbf{E},\mathbf{F}_n))\) 和 \((\mathrm{Hom}_{\mathbf{A}}(\mathbf{A}^r,\mathbf{F}_n))\) 相容。由于 \(\mathrm{Hom}_{\mathbf{A}}(\mathbf{E},\mathbf{F})\) 和 \(\mathrm{Hom}_{\mathbf{A}}(\mathbf{A}^r,\mathbf{F})\) 是有限生成的（引理 2），根据命题 1，只需证明滤过 \((\mathrm{Hom}_{\mathbf{A}}(\mathbf{A}^{r},\mathbf{F}_{n}))\) 是 m-好的；但这由典范同构 \(\mathrm{Hom}_{\mathbf{A}}(\mathbf{A}^{r},\mathbf{F}_{n})\to \mathbf{F}_{n}^{r}\) 的存在以及关系 \(\mathfrak{mF}_n = \mathbf{F}_{n + 1}\) 蕴含 \(\mathfrak{m}(\mathbf{F}_n^r) = (\mathfrak{m}\mathbf{F}_n)^r = \mathbf{F}_{n + 1}^r\) 这一事实立即得到（代数，第二章，第 3 节，第 7 号，注）。

命题 3。设 A 为 Noether 环，m 为 A 的理想，且 A 关于 m-进拓扑是 Hausdorff 且完备的。设 E 是滤过环 A 上的滤过 A-模，E 的滤过（\(E_{n}\)）满足 \(E_{0} = E\)，且 E 关于由（\(E_{n}\)）定义的拓扑是 Hausdorff 的。则下列条件等价：

\begin{enumerate}
\def\labelenumi{(\alph{enumi})}
\item
  E 是有限生成的 A-模，且 \((\mathbf{E}_{n})\) 是 m-好滤过。
\item
  gr(E) 是有限生成的 gr(A)-模。
\item
  对所有 \(n \geqslant 0\)，\(\mathrm{gr}_n(\mathbf{E})\) 都是有限生成的 A-模，并且存在 \(n_0\)，使得当 \(n \geqslant n_0\) 时典范同态

\[
\operatorname{gr} _ {1} (\mathrm{A}) \otimes_ {\mathrm{A}} \operatorname{gr} _ {n} (\mathrm{E}) \rightarrow \operatorname{gr} _ {n + 1} (\mathrm{E})\tag{3}
\]

是满射。
\end{enumerate}

由定义立即可见，(a) 蕴含 (c)。(b) 蕴含 (c) 是第 1 节第 3 号引理 1 的结果；反之，若 (c) 成立，则显然 gr(E) 作为 gr(A)-模由各 gr \(_{p}\) (E)（p ≤ n \(_{0}\)）的和生成，因而根据假设有有限生成系。还需证明 (c) 蕴含 (a)；由于各 gr \(_{n}\) (E) 都是有限生成的且 E \(_{0}\) = E，首先显然可由 n 归纳得出，对所有 n，E/E \(_{n}\) 是有限生成的 A-模；因此只需证明，当 n \textgreater{} n \(_{0}\) 时，E \(_{n}\) 是有限生成的 A-模且 mE \(_{n}\) = E \(_{n+1}\)。现在，考虑 A-模 E \(_{n+1}\)，赋予它穷竭且分离的滤过 (E \(_{n+k}\) )(k ≥ 1)；mE \(_{n}\) ⊂ E \(_{n+1}\)；假设 (c) 蕴含，mE \(_{n}\) 在 gr \(_{n+1}\) (E) = E \(_{n+1}\) /E \(_{n+2}\) 中的像等于 gr \(_{n+1}\) (E)，并生成分次 gr(A)-模 gr(E \(_{n+1}\))。由于按假设 gr \(_{n+1}\) (E) 是有限生成的 A-模，故由第 2 节第 9 号命题 12 得 mE \(_{n}\) = E \(_{n+1}\)，且 E \(_{n+1}\) 是有限生成的 A-模。

